\begin{quote}\small Historical Round26 roadmap; subsequent integrated chapters record its execution.\end{quote}
\section{Next planned goals and remaining proof obligations}
\label{sec:roadmap}

Round26 completes AB1/AB2, AC1/AC2, AD1/AD2, AE1/AE2, and AF1/AF2. The first three goal pairs were inherited from the earlier roadmap. AE and AF were selected only after those six reviews. This paper does not execute further physics loops. The next three goals below were planned from the final evidence. Their second loops must be chosen after their first-loop reviews; two later goals remain unselected.

\begin{enumerate}
\item \textbf{AG: the first actual nonlinear homogeneous update.}
Start within AE2's restricted coupling region. Derive the generated source, scalar term, and diagonal update from one actual transformation. Bound every translated crossing and the relevant unbounded domains in the original rooted support-cardinality norm. The decisive test is whether per-source improvement also improves the collective weighted budget. A failure would select a revised construction, not an assertion of all-stage induction.
\item \textbf{AH: larger physical graphs with complete channel control.}
Freeze the next graph and build its actual Gauss-invariant retained basis. Establish all Gram, coupling, and omitted-channel identities before comparing accuracy. Compute the graph dependence of residual, ground-projection, and true-output denominator constants. A favorable graph-size trend remains weaker than a volume-uniform theorem.
\item \textbf{AI: observable and physical-scale identification.}
Freeze actual Wilson observables, independent parameters, and physical observation times. Determine the parameter combinations identified by their complete readouts, using null probes and indistinguishable parameter families as controls. Add a genuinely discriminating observable or a state/domain/generator map with a bounded defect. A fitted clock that forces agreement is not an admissible substitute.
\end{enumerate}

Beyond these near-term goals, a continuum implication requires an explicitly matched scale trajectory, control uniform in graph volume and cutoff, a nontrivial limiting theory, the required reconstruction/positivity properties, and a positive gap in physical units. These are separate obligations, not a consequence of the finite-graph error threshold. The present results provide controlled subproblems and counterexamples that organize this work; they do not settle it.

